\begin{definition}[The round subsystem]
  \label{def:aba-gbcasub}
  \lean{PLTS.ABA.GSub.GEvt}\lean{PLTS.ABA.GSub.GLab}\lean{PLTS.ABA.GSub.gEvents}\lean{PLTS.ABA.GSub.GProcStep}\lean{PLTS.ABA.GSub.gbcaProc}\lean{PLTS.ABA.GSub.GNetState}\lean{PLTS.ABA.GSub.GNetState.corrupt}\lean{PLTS.ABA.GSub.GNetStep}\lean{PLTS.ABA.GSub.gNet}\lean{PLTS.ABA.GSub.GSubState}\lean{PLTS.ABA.GSub.subPre}\lean{PLTS.ABA.GSub.sub}\lean{PLTS.ABA.GSub.gOwns}\lean{PLTS.ABA.GSub.gAct}\lean{PLTS.ABA.GSub.gbcaSide}\leanok
  \uses{def:aba-procnet, def:aba-netadv, def:aba-gbca-vocab, def:relabel, def:syncproduct, def:idle-family, def:process-algebra}
  One round of graded agreement, taken apart into the pieces that run it: $n$
  local stage programs beside that round's own message fabric.

  The local program $\mathrm{gbcaProc}(P,r,j)$ holds one stage node
  $\mathrm{ProcNodeN}(n)$ of Definition~\ref{def:aba-procnet}: process $j$'s
  stage record for round $r$ and its inbox rows, indexed by sender.
  It holds no corrupted set, no corruption flag and no record of what it has
  multicast.
  Its twenty-five rows are the stage-visible half of the rows of
  Definition~\ref{def:aba-procnet} that touch round $r$: the call, the eight
  multicasts of the message ladder, the delivery, the three graded returns,
  the five Byzantine drive rows, and the rows on which the record does not
  move, either because another process is named or because the call is one
  against an already-opened record.
  The round loop is not here.
  A call writes the stage record alone, and a return sets that record's
  $returned$ flag alone.

  The round's message fabric $\mathrm{gNet}(P,r)$ holds the per-sender pools
  $pool(j)$ of the stage messages $j$ has multicast in this round (D5) and the
  corrupted set $F$.
  Its nine rows are the network half of the same rules: pooling on a multicast,
  the membership guard $m \in pool(j)$ on a delivery, the pooling of the
  $\langle\mathrm{INPUT},b\rangle$ that a call multicasts, the Byzantine
  injection $\mathrm{byzG}$, and the rows on which nothing is sent.
  Corruption is not a row of either table: it is the act the family of
  Definition~\ref{def:idle-family} applies to every round's fabric at once,
  inserting $k$ into $F$ when $k \notin F$ and $|F| < f$ (D1).

  The two rendezvous of one round are the multicast $\mathrm{snd}(j,m)$ and the
  delivery $\mathrm{dlv}(i,j,m)$, which make up the alphabet
  $\mathrm{GEvt}(n)$.
  They carry no round tag: the round is the identity of the instance they
  belong to.
  The subsystem is
  \[
    \mathrm{subPre}(P,r) \;=\;
      \bigl(\mathrm{syncProduct}_{j \in [n]}\ \mathrm{gbcaProc}(P,r,j)\bigr)
        \parallel \mathrm{gNet}(P,r),
  \]
  \[
    \mathrm{sub}(P,r) \;=\; \mathrm{relabel}\bigl(
      \mathrm{abstract}_{\mathrm{GEvt}(n)}(\mathrm{subPre}(P,r))\bigr),
  \]
  over the subsystem-internal alphabet
  $\mathrm{GLab}(n) = \mathrm{NLab}(n) \oplus \mathrm{GEvt}(n)$ and then over
  $\mathrm{NLab}(n)$.
  The two rendezvous are hidden before anything outside the subsystem sees it,
  so the interface of $\mathrm{sub}(P,r)$ is the round's ports and nothing
  else: $\mathrm{callG}(r,\cdot,\cdot)$, $\mathrm{retG}(r,\cdot,\cdot)$,
  $\mathrm{gcallLoop}(r,\cdot,\cdot)$ and the three Byzantine
  graded-agreement drives of round $r$.
  The stage multicast $\mathrm{gsnd}$ and the stage delivery $\mathrm{gdlv}$ of
  Definition~\ref{def:aba-procnet} are labels no component of the subsystem
  offers, so they carry no transition of it.

  A Byzantine drive is split between two components (D11).
  The subsystem carries its effect and not its authorisation: a driven call
  opens the stage record and pools its $\langle\mathrm{INPUT},b\rangle$ under no
  $k \in F$ guard, and a driven return sets $returned$ on the same SEAL-level
  evidence an honest return needs.
  A drive label therefore stays visible at this boundary, and the guard
  $k \in F$ is read outside.

  The graded-agreement side of the deployed protocol is the
  $\mathbb{N}$-indexed family
  \[
    \mathrm{gbcaSide}(P) \;=\;
      \mathrm{family}\bigl(\mathrm{sub}(P,\cdot),\ \mathrm{gOwns},\
        \mathrm{fail},\ \mathrm{gAct}(P)\bigr),
  \]
  where $\mathrm{gOwns}$ sends each of the six interface label classes to its
  round tag and every other label of $\mathrm{NLab}(n)$ to no round.
  A round-tagged label moves its round alone, $\tau$ moves one round, and
  $\mathrm{fail}(k)$ is the broadcast that keeps every round's copy of the
  corrupted set in lockstep.
\end{definition}
